%% =====================================================================
%%  T-15-02  Uncertainty and Similarity
%%  -------------------------------------------------------------------
%%  One idea per file. Core is mandatory. Slots in canonical order.
%%  Max 700 words. No layout commands. No filler. New source material
%%  for this entry goes in sources/<BOOK>/T-15-02.tex -- never by
%%  rewriting this file (docs/RULES.md R0.1).
%% =====================================================================
%% @kind: topic
%% @id: T-15-02
%% @title: Uncertainty and Similarity
%% @subtitle: The two conditions that switch the crowd on
%% @chapter: c15
%% @order: 20
%% @status: stable
%% @sources: B4
%% @updated: 2026-09-19

\topic{T-15-02}{Uncertainty and Similarity}{The two conditions that switch the crowd on}

\begin{core}
Crowd evidence is not consulted evenly. It is followed hardest in two conditions: when the situation is ambiguous, and when the crowd is made of people like you. Either one raises compliance; together they raise it most.
\end{core}
\begin{mechanism}
Ambiguity removes the alternative. Where the facts are clear a person can check them; where they are not, there is nothing to check against and other people's behaviour becomes the only available data. Urgency compounds this, because it removes the time in which ambiguity could be resolved --- which is why a crowd is most influential in exactly the situations where it is least reliable.

Similarity supplies the relevance test. Other people's behaviour is informative about what someone like you should do, and uninformative about what anyone should do. A crowd of strangers is weak evidence; a crowd of peers is strong. This is why the technique is worked by manufacturing resemblance --- accent, age, dress, occupation, background --- rather than by manufacturing numbers alone, and why a single well-matched witness can outweigh a large unmatched crowd.
\end{mechanism}
\begin{tells}
  \item You cannot state the facts of the situation, only what others seem to think they are.
  \item A decision is required quickly and information is arriving slowly.
  \item The people being cited resemble you specifically --- your age, your role, your circumstances.
  \item The resemblance is mentioned rather than incidental.
  \item Nobody is investigating; everyone is watching everyone else.
\end{tells}
\begin{moves}
  \item Create ambiguity before presenting the crowd: withhold the facts that would let the counterpart check independently.
  \item Add time pressure so the ambiguity cannot be resolved by enquiry.
  \item Match the witnesses to the counterpart. One similar person is worth several dissimilar ones.
  \item Lead with resemblance and follow with numbers, in that order.
\end{moves}
\begin{counters}
  \item Reduce ambiguity first and consult the crowd second. Get the facts from a source the other party does not control, then decide.
  \item Buy time deliberately. Almost every crowd-driven decision improves overnight, and the pressure to decide now is itself a signal.
  \item Discount resemblance that is presented to you. A witness selected because they are like you was selected for that reason.
  \item Ask what the dissimilar people are doing. If only peers agree, the agreement may be a property of the peer group rather than of the question.
\end{counters}
\begin{cost}
Both conditions are also the normal conditions of ordinary life. Refusing to consult others when uncertain, or refusing to weight peer experience, produces a person who reinvents every judgement from nothing and gets many of them wrong. The usable rule is narrow: treat crowd evidence as weak exactly where the stakes are high and the ambiguity is manufactured --- and notice that those two tend to arrive together.
\end{cost}

\sources{B4}
\seesources
\seealso{T-15-01, T-13-01, T-14-01, T-23-01}
